\reading{IM-16}{Liquidity Risk Management}{STRIKE FIRST}{2}

\srcnote{Bennett W. Golub (ed.), \textit{BlackRock's Guide to Fixed-Income Risk
Management} (Wiley, 2024), Chapter 7. No GARP source book covers this reading.}

\begin{imtrapbox}[breakable=false,title={The title carries no parenthetical}]
GARP's title for this reading is simply \term{Liquidity Risk Management}. The
``(investment funds)'' the transcribed spine carried was an editorial addition,
not GARP's wording.
\end{imtrapbox}

\begin{imkeybox}[title={The five objectives in this reading}]
\begin{enumerate}[label=\textbf{\alph*.}]
\item Explain best practices for managing liquidity risk in an investment fund.
\item Explain potential challenges fund managers face when managing liquidity risk in fixed-income portfolios, including challenges related to data modeling.
\item Describe and apply different approaches to model asset liquidity, including an approach to modeling infrequently traded bonds and the t-cost model to calculate transaction costs for corporate bonds.
\item Identify examples of extraordinary measures that regulators may allow funds to take to meet unanticipated redemption requests, and explain when each measure may be allowed.
\item Explain approaches to manage and model redemption risk in an investment fund, including the redemption waterfall, approaches to model a fund's level of redemption-at-risk, and approaches to optimize liquidity to meet redemption requests.
\end{enumerate}
\end{imkeybox}

% ------------------------------------------------------------------
\lo{IM-16 a}{Explain best practices for managing liquidity risk in an investment fund.}

\begin{imdefbox}[title={Fund liquidity risk}]
The risk that a fund may be \term{unable to obtain enough liquidity to meet
investor redemptions}. It can arise as a \emph{byproduct} of other issues, such as
poor fund performance or reputational problems involving managers.
\end{imdefbox}

The most effective response is a robust, \term{top-down governance structure}
supported by a liquidity risk management framework built on four elements:
\term{liquidity risk measurement} (identify and quantify risks; assess how quickly
assets can be sold); \term{liquidity risk management} (monitor the liquidity
profile continuously; compare redemption risk against limits); \term{portfolio
construction}; and \term{governance}.

\begin{imkeybox}[title={The core objective has two sides}]
Ensure \term{enough cash to meet redemptions} --- \emph{but} avoid
\term{leaving remaining investors with the illiquid assets}. Selling the liquid
holdings first meets today's redemption and hands tomorrow's problem to whoever
stayed.
\end{imkeybox}

% ------------------------------------------------------------------
\lo{IM-16 b}{Explain potential challenges fund managers face when managing liquidity risk in fixed-income portfolios, including challenges related to data modeling.}

Following the Global Financial Crisis, low interest rates led to a surge in
corporate borrowing, including more complex bond issuances --- and a shift in
market behaviour.

\begin{imkeybox}[title={Price takers to price makers}]
\term{Before the crisis}, asset managers were \term{price takers}, relying on
dealer quotes. \term{After the crisis} they became \term{price makers}, setting
prices and searching for counterparties. The consequences: \term{thinner markets}
with lower trading volume, \term{wider bid-ask spreads}, and \term{higher price
impact} for large trades.
\end{imkeybox}

Large trades are often split into smaller orders, but liquidity becomes more
expensive when trades must be executed quickly, especially for complex securities.

\subsection{Data availability}

\renewcommand{\arraystretch}{1.2}
\noindent\begin{tabularx}{\textwidth}{@{}>{\raggedright\arraybackslash}p{2.6cm}X@{}}
\toprule
TRACE, as it was & FINRA's TRACE system provides bond price and volume data, but \term{trade sizes are capped}: high-yield trades above \term{\$1 million} are reported as ``\$1M+'' and investment-grade trades above \term{\$5 million} as ``\$5M+''. This creates ambiguity --- \term{\$6 million and \$60 million look the same} --- making liquidity estimation less accurate. Models try to infer actual volumes from historical patterns, but this can fail in changing market conditions. \\
TRACE now & Provides \term{uncapped} trade volumes at CUSIP level, but with a \term{six-month delay}. \\
MiFID II & Intraday price and volume data for \emph{liquid} bonds; \term{one-month delayed} data for less liquid bonds; and \term{mandatory reporting of transaction costs}. \\
Proposed & A \term{consolidated tape} in Europe, combining data across venues. \\
\bottomrule
\end{tabularx}
\renewcommand{\arraystretch}{1}
\figcap{The data problem in one table. Note what the fixes trade away: the cap is
removed only at the price of a six-month delay, so the data becomes accurate
exactly when it stops being timely. The two numbers most worth carrying are the
\term{\$1 million high-yield} and \term{\$5 million investment-grade} caps ---
the higher threshold applies to the \emph{more} liquid category.}

% ------------------------------------------------------------------
\lo{IM-16 c}{Describe and apply different approaches to model asset liquidity, including an approach to modeling infrequently traded bonds and the t-cost model to calculate transaction costs for corporate bonds.}

\subsection{The days-to-liquidate approach}

Estimates how long it takes to unwind a position, from three components:

\begin{imfmlbox}
\[ \text{days to liquidate} = \frac{\text{trade size}}{\text{participation rate} \times \text{ADV}} \]
\vk{\term{trade size} is based on redemptions, leverage and derivatives exposure;
\term{participation rate} is the percentage of average daily volume that can be
sold without major price impact --- hard to estimate, and related to \term{latent
liquidity}; \term{ADV} is average trading volume over a period, estimable from
market data or proxies.}
\end{imfmlbox}

\subsection{Infrequently traded bonds}

When bonds do not trade frequently, direct estimation becomes unreliable. A
\term{random forest regression} --- a machine learning technique --- is well suited
to \term{nonlinearity, missing data and large sets of explanatory variables},
assuming expected volume on no-trade days is zero. Analysts supplement the data
using \term{comparable bonds} (similar risk and characteristics) and \term{dealer
inventory data} (bid-ask quotes indicating potential supply and demand).

\begin{imtrapbox}[breakable=false,title={Quoted liquidity is not executed liquidity}]
Both supplementary sources have the same limitation: \term{quoted liquidity is
often higher than actual executed liquidity}, which points back to latent
liquidity. To address the uncertainty, liquidity is often expressed as a
\term{range rather than a single number}, estimated using tail measures. Note also
the direction on volume variability: when variance is high, \term{ADV may
understate potential liquidity}, making it a \emph{conservative} measure.
\end{imtrapbox}

\subsection{Transaction cost (t-cost) modelling}

T-cost modelling assesses how costly it is to liquidate positions, balancing two
competing objectives:

\begin{imkeybox}
\term{Liquidate quickly} $\rightarrow$ higher price impact. \quad
\term{Maintain liquidity} $\rightarrow$ lower returns or delayed execution.
\end{imkeybox}

It uses intraday benchmark prices, empirical transaction data and bond-specific
attributes to forecast both \term{fixed costs} and \term{market impact costs},
with regression techniques estimating certain components.

\begin{imdefbox}[title={Implementation shortfall}]
The primary purpose of t-cost modelling is to estimate \term{market impact},
measured as the difference between the \term{traded price} (actual execution) and
the \term{order-arrival benchmark price} or another appropriate pre-trade
benchmark.
\end{imdefbox}

\begin{imtrapbox}[breakable=false,title={Which benchmark, and how to weight}]
Using the \emph{previous end-of-day price} as the benchmark is misleading, because
it includes market-driven price movements that occurred \emph{before} execution.
An \term{intraday pre-trade benchmark} is best practice. When estimating costs,
\term{weight observations by trade size} to avoid bias from many small trades, and
use security-level data where available. Machine learning methods --- clustering,
tree models, neural networks --- further improve estimates.
\end{imtrapbox}

% ------------------------------------------------------------------
\lo{IM-16 d}{Identify examples of extraordinary measures that regulators may allow funds to take to meet unanticipated redemption requests, and explain when each measure may be allowed.}

Advance planning remains the most effective way to manage redemption pressure, but
funds may need extraordinary measures when normal liquidity tools are
insufficient. These are \term{rarely used but must be deployable quickly}.

\begin{imkeybox}[title={``Break glass'' testing}]
Legal, operations, risk management and trading teams should conduct
\term{``break glass'' testing} to ensure each measure can be activated quickly if
required. A tool that exists on paper but has never been exercised is not a tool.
\end{imkeybox}

\term{Temporary borrowing} provides a short-term liquidity buffer, allowing the
fund to meet redemptions without immediately selling illiquid assets. Common
sources include \term{repos, overdrafts and credit facilities}. It avoids forced
selling, although it is permitted only within regulatory limits.

% ------------------------------------------------------------------
\lo{IM-16 e}{Explain approaches to manage and model redemption risk in an investment fund, including the redemption waterfall, approaches to model a fund's level of redemption-at-risk, and approaches to optimize liquidity to meet redemption requests.}

Outflows can be driven by macroeconomic developments, poor fund performance,
reputational issues, or fund-specific factors such as investor profiles and
concentration.

\begin{imkeybox}[title={The main objective}]
Manage redemptions \term{fairly and equitably across all investors} --- those
redeeming, those subscribing, and \emph{those remaining in the fund}. The third
group is the one an ad hoc response forgets.
\end{imkeybox}

The \term{redemption waterfall} for '40 Act funds is a multi-layer approach to
meeting redemption requests when assets are insufficient at one level. Stress
testing is a natural extension of redemption modelling.

\subsection{Modelling redemption-at-risk}

Analysts estimate worst-case outflows --- for example a \term{99th-percentile,
five-day scenario} --- using historical flows. This forms the \term{HRaR model}
and provides stress signals. Three limitations:

\begin{enumerate}[label=\alph*)]
\item \term{Limited history.} New or fast-growing funds may not show meaningful
outflows.
\item \term{Unpredictable behaviour.} Investor actions can be idiosyncratic.
\item \term{Data gaps.} Managers may lack visibility due to intermediaries such as
custodians.
\end{enumerate}

Large outflows usually result from a mix of \term{systemic factors} (macro shocks,
sector-wide disruptions) and \term{idiosyncratic factors} (investor behaviour, peer
performance, cash needs, reputational or regulatory events). Machine learning can
improve HRaR by capturing multiple factors and nonlinearities, helping estimate an
\term{upper bound} of redemptions.

\subsection{Liquidity optimization}

Fund liquidity management can be framed as an \term{optimization problem} in which
managers balance liquidity, transaction costs and portfolio risk while meeting
redemption obligations. Each security $i$ has a market value $m_i$, a notional
value $n_i$, and a proportion sold $p_{i,t}$; the first step is to determine the
optimal \term{pro rata} fraction, and the objective is then maximised subject to
constraints.

\renewcommand{\arraystretch}{1.2}
\noindent\begin{tabularx}{\textwidth}{@{}>{\raggedright\arraybackslash}p{2.6cm}X@{}}
\toprule
\textbf{\sffamily\small Constraint} & \textbf{\sffamily\small What it enforces} \\
\midrule
Redemption & Total liquidation meets the required outflow $v$ over a specified period --- for example \$200 million in two days. Guarantees sufficient cash generation. \\
ADV & Security $i$ has an ADV of $v_i$ and belongs to sector $j$ with an assigned \term{participation rate} $r_j$; sales are capped accordingly. \\
Time-lag & Some assets cannot be sold immediately ($d_i$ days of delay --- IPO lockups, for instance). Reflects real-world restrictions. \\
Risk & Well-managed portfolios also stay \term{below a specified risk threshold} $R$ after the sales. \\
\bottomrule
\end{tabularx}
\renewcommand{\arraystretch}{1}
\figcap{Four constraints, and a warning attached to all of them: \term{the
optimization is ultimately constrained by the quality and availability of the
model's input data.} Every constraint above needs a number that IM-16 b has just
finished explaining is hard to observe.}

\begin{imtrapbox}[title={Consolidated trap summary --- IM-16}]
\begin{itemize}
\item \term{Definition, IM-16 a.} Fund liquidity risk is about meeting
\emph{redemptions}, and it arises as a \emph{byproduct} of performance or
reputational problems.
\item \term{Polarity, IM-16 b.} Managers moved from price \emph{takers} to price
\emph{makers} after the crisis, and the result was \emph{worse} liquidity: thinner
markets, wider spreads, higher impact.
\item \term{Number, IM-16 b.} TRACE caps are \$1 million for \emph{high-yield} and
\$5 million for \emph{investment-grade}. The higher cap applies to the more liquid
category. Uncapped data now exists but with a six-month delay; MiFID II's delay for
illiquid bonds is one month.
\item \term{Polarity, IM-16 c.} \term{Quoted liquidity exceeds executed
liquidity}, and high volume variance means ADV \emph{understates} liquidity ---
making it conservative, not aggressive.
\item \term{Definition, IM-16 c.} Implementation shortfall uses the
\term{order-arrival} or intraday pre-trade benchmark, \emph{not} the previous
close.
\item \term{Role, IM-16 c.} Weight observations by \emph{trade size}; unweighted
data is biased by many small trades.
\item \term{Scope, IM-16 e.} Fairness runs to three groups: redeeming,
subscribing, and \emph{remaining}.
\end{itemize}
\end{imtrapbox}

% ==================================================================
